\chapter{Conclusion}

While the model audited in this project demonstrates strong predictive accuracy, we find that the underlying dataset poses significant challenges for responsible deployment in an automated decision system (ADS). The data is synthetically generated with unclear provenance and shows substantial geographic and demographic biases, particularly overrepresentation of urban centers in North and Central-West India and underrepresentation of rural, southern, and northeastern regions. Moreover, culturally relevant variables such as family structure, interpersonal relationships, and stigma-related experiences, which are central to understanding depression in the Indian context, are missing. As a result, while the dataset is useful for experimentation and benchmarking, it is not appropriate for real-world, fairness-sensitive applications without considerable expansion and refinement.

The implementation itself is robust in terms of accuracy, with the final ensemble achieving over 94\% on held-out validation data. However, fairness evaluations tell a different story. While gender-based disparities were minimal, significant disparities emerged along employment status: students were over-identified as depressed (high false positive rate), while working professionals were often missed (high false negative rate). These findings were confirmed through \texttt{MetricFrame} audits, confusion matrices, and local explanations using \texttt{LIME}. To address these disparities, we applied \texttt{Fairlearn}’s \texttt{ThresholdOptimizer} with an equalized odds constraint. This post-processing method successfully reduced group-level disparities in true and false positive rates, at the cost of a modest drop in accuracy ($\sim 4\%$). Given the sensitive nature of mental health prediction, we believe this is an acceptable trade-off. Stakeholders such as clinicians, patients, and policy makers would likely find fairness metrics, particularly false negative and false positive rates, more meaningful than accuracy alone, as these capture the real-world harms of misclassification.

Despite improvements, we would not recommend deploying this model in the public sector or industry in its current form. The lack of transparency in data sourcing, insufficient cultural representation, and observed subgroup disparities make it unsuitable for high-stakes decision-making without further development. Future iterations should incorporate richer, more representative data; include culturally and socially relevant features; apply fairness interventions during model training rather than solely post hoc; and calibrate prediction outputs to support informed threshold selection. Only with these enhancements can such a model be responsibly used to support equitable and trustworthy mental health screening at scale.